\documentclass[11pt,a4paper]{article}

% ---------- Page layout ----------
\usepackage[margin=1in]{geometry}
\usepackage{setspace}
\setstretch{1.1}
\usepackage[T1]{fontenc}
\usepackage{lmodern}

% ---------- Section formatting ----------
\usepackage{titlesec}
\titleformat{\section}{\large\bfseries}{}{0em}{}
\titlespacing*{\section}{0pt}{1.2em}{0.6em}

% ---------- Bibliography ----------
% \usepackage[numbers]{natbib}
% \bibliographystyle{unsrt}

% \usepackage[backend=biber,style=numeric,sorting=none]{biblatex}
% \addbibresource{publications.bib} % your .bib file

% \usepackage{bibentry}
% \PassOptionsToPackage{style=ieee,sorting=ydnt,uniquename=init,defernumbers=true}{biblatex}
% % Remove default [Online] prefix before URLs
% \DefineBibliographyStrings{english}{url={\textsc{url}}}
% % Add short journal field formatting after loading settings
% \DeclareFieldFormat{journaltitle}{%
%   \iffieldundef{shortjournal}{#1}{\printfield{shortjournal}}}
% \addbibresource{publications.bib}

\PassOptionsToPackage{
    style=ieee,
    sorting=ydnt,
    uniquename=false,
    defernumbers=true,
    maxnames=99,
    minnames=99,
    doi=true,
    url=false,
    eprint=false
}{biblatex}
\usepackage[backend=biber,style=numeric]{biblatex}

% \usepackage{biblatex}
% Shorten "url" prefix for IEEE style
\DefineBibliographyStrings{english}{url={\textsc{url}}}
% Load your bibliography
\addbibresource{publications.bib}

% ---------- No page numbers ----------
% \pagestyle{empty}

% add hlines under the section titles
\usepackage{titlesec}
\titleformat{\section}
  {\large\bfseries}
  {}
  {0em}
  {}
  [\titlerule]

% links
\usepackage[hidelinks]{hyperref}

\usepackage{enumitem}

\begin{document}


% ---------- Header ----------
\begin{center}
{\Large \textbf{Yuzhe Zhang}} \\[0.3em]
\today{} \\
\end{center}

\vspace{0.5em}

% ---------- Sections ----------
\section*{Personal Information}

\noindent\textbf{Date and place of birth:} January 10, 1998 (Jiangsu, China) \\
\textbf{Nationality:} Chinese \\
\textbf{Address:} GSI Helmholtz Institute Mainz, Staudingerweg 18, 55128 Mainz, Germany \\
\textbf{Email:} yuhzhang@uni-mainz.de

\noindent\textbf{ORCID:} \href{https://orcid.org/0009-0006-6082-9066}{\texttt{0009-0006-6082-9066}}

\noindent\textbf{Google Scholar:} \href{https://scholar.google.com/citations?user=tPAwEPsAAAAJ&hl=en&oi=sra}{scholar.google.com/citations?user=tPAwEPsAAAAJ}

\noindent\textbf{LinkedIn:} \href{https://www.linkedin.com/in/yuzhe-zhang-physics/}{linkedin.com/in/yuzhe-zhang-physics/}

\noindent\textbf{Webpage:} \href{https://yuzhe98.github.io/}{yuzhe98.github.io}

\noindent\textbf{Languages:} Chinese (native), English (professional working proficiency)

\section*{Summary}
I am a junior researcher focusing on \textbf{axionlike dark matter} detection using \textbf{nuclear magnetic resonance} and \textbf{quantum sensing} techniques. 
Within the Cosmic Axion Spin Precession Experiment (CASPEr) collaboration, I have contributed to experiments, data analysis, and simulations, resulting in peer-reviewed publications and presentations at international conferences and seminars. 
Over the years, I have developed expertise in precision measurement, programming, and electronics. 
% Beyond axion, I am \textbf{broadly interested in dark matter} research. 
I also prioritize collaborative research, clear scientific communication, and teaching.

\section*{Education}

\textbf{\textit{Ph.D.}, Physics (in progress)}\hfill 11.2022--Present \\
Johannes Gutenberg University Mainz \\
Thesis: \emph{Search for an Axionlike Dark Matter Candidate Using Nuclear Magnetic Resonance}

\vspace{0.5em}

\noindent\textbf{Master of Science, Physics} \hfill 10.2020--10.2022 \\
Johannes Gutenberg University Mainz\\
Thesis: \emph{Search for axion-like particles by performing NMR on a thermally polarized sample}

\vspace{0.5em}

\noindent\textbf{Bachelor of Science, Physics} \hfill 09.2016--06.2020 \\
University of Science and Technology of China\\
Thesis: \emph{Data analysis in Earth rotation measurements based on a nuclear spin--electron spin system}


\section*{Research}
\noindent\textbf{Cosmic Axion Spin Precession Experiment (CASPEr)} \hfill 10.2020--Present \\
GSI Helmholtz Institute Mainz \& Johannes Gutenberg University Mainz \\
Advisor: Prof.\,Dr.\,Dmitry Budker\\
\textbf{Contribution:} Experiment, simulation, and data analysis. Design and optimization of axionlike dark matter detection strategy. Cross-team coordination within the CASPEr collaboration. 

\vspace{0.5em}

\noindent\textbf{Rb--Xe Comagnetometer Experiment} \hfill 09.2019--06.2020 \\
Zheng-Tian Lu Group, University of Science and Technology of China \\
Advisor: Prof.\,Dr.\,Dong Sheng\\
\textbf{Contribution:} Assistance in the experimental search for spin-gravity coupling. Analysis of co-magnetometer data. 

\vspace{0.5em}

\noindent\textbf{Global Network of Optical Magnetometers for Exotic physics (GNOME)} \hfill 06.2019--08.2019 \\
Laboratory of Magnetometry and Quantum-State Engineering, Jagiellonian University \\
Advisor: Prof.\,Dr.\,Szymon Pustelny\\
\textbf{Contribution:} Simulation and data analysis for the GNOME collaboration.

\vspace{0.5em}

\noindent\textbf{Catalyst Research for Water Splitting} \hfill 08.2017--06.2019 \\
Jie Zeng Group, University of Science and Technology of China \\
Advisor: Prof.\,Dr.\,Jie Zeng\\
\textbf{Contribution:} Synthesized catalysts, prepared electrodes, and tested water-splitting performance.


\section*{Selected Publications}

\noindent

\AtNextCite{\defcounter{maxnames}{3}\defcounter{minnames}{3}}
\noindent[1] \fullcite{walter2025searchaxionlikedarkmatter}
\begin{itemize}
    \item Commissioning results of the CASPEr-Gradient experiment.
    \item Co--first author (equal contribution; 3 authors). Contributions: apparatus construction and commissioning; data acquisition; analytical calculations and data analysis; manuscript preparation; coordination within the collaboration.
\end{itemize}

\AtNextCite{\defcounter{maxnames}{3}\defcounter{minnames}{3}}
\noindent[2] \fullcite{Yuzhe2023_FrequencyScanning}
\begin{itemize}
    \item Theoretical considerations on frequency scanning for axion haloscopes. 
    \item First and corresponding author. Contributions: all theoretical calculations and analysis; full manuscript preparation; coordination within the collaboration.
\end{itemize}

% \AtNextCite{\defcounter{maxnames}{3}\defcounter{minnames}{3}}
% \noindent[3] \fullcite{Walter2023Shimming}
% \begin{itemize}
%     \item Application of Bayesian optimization to accelerate magnetic field shimming.
%     \item Contributions: experimental apparatus construction; data acquisition; optimization of the algorithm.
% \end{itemize}

\section*{Presentations}
\noindent\textbf{Talks}

\vspace{0.3em}

\noindent\begin{tabular}{@{}lp{14cm}@{}}
2026 & GNOME workshop \\
2023 & Academic Seminar of Chinese Students and Scholars in the German--French Region \\
2022 & Invited talk at a Seminar at Institute of Theoretical Physics, Chinese Academy of Sciences \\
2021 & 16th Patras Workshop on Axions, WIMPs and WISPs \\
\end{tabular}

\vspace{0.5em}

\noindent\textbf{Posters}

\vspace{0.3em}

\noindent\begin{tabular}{@{}lp{14cm}@{}}
2024 & The Fundamentals School -- summer school in high-energy and gravitational physics \\
2024 & DPG Spring Meeting of the Atomic, Molecular, Quantum Optics and Photonics Section (SAMOP) \\
2023 & The conference on dark matter, ions, molecules and atoms \\
2022 & 17th Patras Workshop on Axions, WIMPs and WISPs \\
\end{tabular}

\section*{Software Development}

\texttt{AxionBloch} -- [\href{https://github.com/Yuzhe98/AxionBloch}{GitHub link}]
\begin{itemize}
    \item Simulation tool for axion-induced spin dynamics using Bloch equations.  
    \item Contributions: all design and implementation. 
\end{itemize}

\section*{Academic / Educational Contributions}

\noindent\begin{tabular}{@{}lp{13.5cm}@{}}
2026 & GNOME workshop  \\
& Local co-organizer \\
2026 & Educational publication – \textit{Some Rules for Good Scientific Presentations} (contribution: manuscript translation and editing)\\
2025--2026 & Teaching assistant in Advanced Physics Laboratory Course - Magnetic Resonance (undergraduate-to-master-level course)\\
2024 & Educational publication – \textit{Some rules of good scientific viewgraphs} (contribution: idea initiation; manuscript writing \& editing)\\
2022 & 17th Patras Workshop on Axions, WIMPs and WISPs  \\
& Local co-organizer
\end{tabular}


\section*{Skills and Experiences}

\noindent\textbf{Experiment / Technical}
\begin{itemize}
    \item \textbf{Nuclear Magnetic Resonance (NMR):} Apparatus construction, operation, calibration, and data acquisition.
    \item \textbf{Quantum sensing / precision measurements:} Operating Superconducting Quantum Interference Device (SQUID).
    \item \textbf{Cryogenics:} Managing liquid helium and nitrogen for experiments.
    \item \textbf{Hardware automation:} Programming-based control of laboratory equipment.
    % \item \textbf{Catalysis / Electrochemistry:} Catalyst synthesis, electrode preparation, and water-splitting performance testing.
\end{itemize}

\noindent\textbf{Programming / Computational}
\begin{itemize}
    \item \textbf{Data analysis:} Processing CASPEr experimental data with Python; knowledge of particle physics data workflows from workshops.
    \item \textbf{Simulation:} Physics simulations (spin dynamics, signal modeling) using C++ and Python.
    \item \textbf{Numerical equation solving:} Solving Bloch and Schrödinger equations using C++ and Python.
    \item \textbf{Software tools:} Experience with Git / Docker / \LaTeX / HDF5.
    \item \textbf{High-performance computing (HPC):} Conceptual understanding and practical usage learned through workshops.
\end{itemize}

\noindent\textbf{Collaboration / Communication}
\begin{itemize}
    \item International teamwork: Collaboration with CASPEr members. Hosting weekly data analysis meetings with international colleagues.
    \item Mentoring and student support: Assisting interns and students with laboratory integration and research projects; teaching assistant in university experiment course.
    \item Scientific communication: Co-authoring publications, poster presentations, and assisting in talks. Organization of international conference. 
\end{itemize}

\section*{Publication List}
\subsection*{Peer-reviewed Journals}
\begin{enumerate}[label={[\arabic*]}]
\item \fullcite{Kim2026haloscopecomparison}
\item \fullcite{walter2025searchaxionlikedarkmatter}
\item \fullcite{Yuzhe2023_FrequencyScanning}
\item \fullcite{Walter2023Shimming}
\item \fullcite{XingYulin2020Electroreduction}
\end{enumerate}

\subsection*{Software}
\begin{enumerate}[label={[\arabic*]}]
% \item \fullcite{Yuzhe2026_axionbloch}
\item \fullcite{Zhang2026_repoAxionBloch}
\end{enumerate}

\subsection*{Preprints}
\begin{enumerate}[label={[\arabic*]}]
\item \fullcite{Yuzhe2026_axionbloch}
% \item \fullcite{Zhang2026_repoAxionBloch}
\end{enumerate}

\subsection*{In Preparation}
\begin{enumerate}[label={[\arabic*]}]
\item \fullcite{Gao2025_EarthDM}
\item \fullcite{Walter2026_CASPErDo}
\end{enumerate}

\end{document}
